\documentclass{article}
\usepackage{amsmath,amssymb,amsthm}
\usepackage[ruled,vlined]{algorithm2e}
\usepackage{enumitem}
\usepackage{hyperref}
\newtheorem{theorem}{Theorem}
\newtheorem{lemma}{Lemma}
\newtheorem{remark}{Remark}
\newcommand{\prox}{\operatorname{prox}}
\newcommand{\grad}{\nabla}
\newcommand{\R}{\mathbb{R}}
\newcommand{\norm}[1]{\left\|#1\right\|}
\begin{document}

\section*{Proximal Gradient Method (Forward–Backward Splitting)}

\section*{Mathematical Formulation}
Let  
\[
F(x)\;:=\;f(x)+g(x),\qquad x\in\R^{d},
\]
where  

\begin{itemize}[nosep]
\item $f:\R^{d}\to\R$ is convex, differentiable and its gradient is $L$‑Lipschitz:
\[
\|\grad f(x)-\grad f(y)\|\le L\|x-y\|\quad\forall x,y\in\R^{d}.
\]
\item $g:\R^{d}\to\R\cup\{+\infty\}$ is proper, closed and convex (possibly non‑smooth).
\end{itemize}

For a stepsize $\eta>0$ the \emph{proximal operator} of $g$ is
\[
\prox_{\eta g}(v)\;:=\;\arg\min_{u\in\R^{d}}\Big\{ g(u)+\frac{1}{2\eta}\|u-v\|^{2}\Big\}.
\]

The algorithm generates a sequence $\{x_{k}\}_{k\ge0}$ by
\begin{equation}
\boxed{\,x_{k+1}= \prox_{\eta g}\bigl(x_{k}-\eta\,\grad f(x_{k})\bigr)\,},\qquad k=0,1,2,\dots
\tag{1}
\end{equation}
The stepsize is required to satisfy $0<\eta\le 1/L$ (see the convergence theorem).

\section*{Pseudocode}
\begin{algorithm}[H]
\caption{Proximal Gradient Method}
\KwIn{Initial point $x_{0}\in\R^{d}$, stepsize $\eta\in(0,1/L]$}
\For{$k\gets0$ \KwTo $T-1$}{
    $v_{k}\gets x_{k}-\eta\,\grad f(x_{k})$\;
    $x_{k+1}\gets\prox_{\eta g}(v_{k})$\;
}
\KwOut{Iterate $x_{T}$}
\end{algorithm}

A common termination criterion is 
\[
\|x_{k+1}-x_{k}\|\le\varepsilon\quad\text{or}\quad
F(x_{k})-F(x_{k+1})\le\varepsilon .
\]

\section*{Dry Run Example}
Consider a one‑dimensional problem
\[
f(w)=(w-3)^{2},\qquad g\equiv0,
\]
so $\prox_{\eta g}$ reduces to the identity.  
Gradient: $\grad f(w)=2(w-3)$.  
Lipschitz constant $L=2$; choose $\eta=0.1\le 1/L=0.5$.

\begin{center}
\begin{tabular}{c|c|c|c}
Iteration $k$ & $w_{k}$ & $\grad f(w_{k})$ & $w_{k+1}=w_{k}-\eta\grad f(w_{k})$\\\hline
0 & $0$ & $2(0-3)=-6$ & $0-0.1(-6)=0.6$\\
1 & $0.6$ & $2(0.6-3)=-4.8$ & $0.6-0.1(-4.8)=1.08$\\
2 & $1.08$ & $2(1.08-3)=-3.84$ & $1.08-0.1(-3.84)=1.464$\\
3 & $1.464$ & $2(1.464-3)=-3.072$ & $1.464-0.1(-3.072)=1.7712$
\end{tabular}
\end{center}

Objective values $F(w_{k})=f(w_{k})$:

\[
f(0)=9,\;f(0.6)=5.76,\;f(1.08)=3.6864,\;f(1.464)=2.3687.
\]

The iterates move monotonically toward the minimiser $w^{\star}=3$.

\newpage
\section*{Assumptions}
\begin{enumerate}[label=\textbf{A\arabic*.},leftmargin=*,itemsep=1ex]
\item \textbf{Convexity and smoothness of $f$.} $f$ is convex and differentiable with $L$‑Lipschitz gradient:
\[
\|\grad f(x)-\grad f(y)\|\le L\|x-y\|,\quad\forall x,y\in\R^{d}.
\tag{A1}
\]
\item \textbf{Convexity of $g$.} $g$ is proper, closed and convex (may be non‑smooth). Its proximal operator is well defined for any $\eta>0$.
\tag{A2}
\item \textbf{Stepsize.} The constant stepsize satisfies
\[
0<\eta\le \frac{1}{L}.
\tag{A3}
\]
\item \textbf{Existence of a minimiser.} The composite problem $\min_{x}F(x)$ admits at least one solution $x^{\star}\in\arg\min F$ and $F(x^{\star})$ is finite.
\tag{A4}
\end{enumerate}

\section*{Main Convergence Theorem}
\begin{theorem}[Sub‑linear convergence of the proximal gradient method]
\label{thm:main}
Under Assumptions \textbf{A1–A4} and with a constant stepsize $\eta\in(0,1/L]$, the iterates generated by \eqref{1} satisfy for every $k\ge0$
\begin{equation}
F(x_{k})-F(x^{\star})\;\le\;
\frac{\norm{x_{0}-x^{\star}}^{2}}{2\eta\,k}.
\tag{2}
\end{equation}
Consequently,
\[
\min_{0\le i\le k-1}\bigl\{F(x_{i})-F(x^{\star})\bigr\}=O\!\left(\frac{1}{k}\right).
\]
\end{theorem}

\section*{Theoretical Properties}
\begin{enumerate}[label=\textbf{P\arabic*.},leftmargin=*,itemsep=1ex]
\item \textbf{Stability via non‑expansiveness.} The proximal operator is firmly non‑expansive:
\[
\|\prox_{\eta g}(u)-\prox_{\eta g}(v)\|^{2}
\le \langle u-v,\;\prox_{\eta g}(u)-\prox_{\eta g}(v)\rangle,
\qquad\forall u,v\in\R^{d}.
\tag{P1}
\]
This property is the cornerstone of the descent inequality used in the proof.
\item \textbf{Stepsize sensitivity.} The bound \eqref{2} holds for any $\eta\le1/L$; a larger $\eta$ (up to the limit $1/L$) yields a smaller constant $1/(2\eta)$ and thus a tighter bound.
\item \textbf{Relation to gradient descent.} If $g\equiv0$, then $\prox_{\eta g}=I$ and \eqref{1} reduces to the classical gradient descent iteration, recovering the standard $O(1/k)$ rate for smooth convex minimisation.
\item \textbf{Comparison to baseline algorithms.}
\begin{itemize}[nosep]
\item Compared with sub‑gradient methods (which need diminishing stepsizes), the proximal gradient method attains the same $O(1/k)$ rate with a \emph{fixed} stepsize.
\item For composite problems where $g$ has a simple proximal map (e.g., $\ell_{1}$ regularisation), the method is often dramatically faster than generic sub‑gradient schemes.
\end{itemize}
\end{enumerate}

\section*{Discussion}
The theorem shows that each iteration yields a guaranteed reduction in the optimality gap that is inversely proportional to the iteration count. The rate is optimal for first‑order methods on the class of convex composite problems with Lipschitz smooth part. The proof hinges on two elementary tools:

\begin{itemize}[nosep]
\item The descent lemma for $L$‑smooth functions.
\item The firm non‑expansiveness (or equivalently, the Moreau decomposition) of the proximal operator.
\end{itemize}
No stochasticity is involved; consequently, expectations do not appear. Extensions to stochastic proximal gradients follow the same template but require variance bounds.

\newpage
\section*{Preliminaries (Definitions and Lemmas)}
\begin{lemma}[Descent Lemma]\label{lem:descent}
If $\grad f$ is $L$‑Lipschitz, then for all $x,y\in\R^{d}$
\[
f(y)\le f(x)+\langle\grad f(x),y-x\rangle+\frac{L}{2}\|y-x\|^{2}.
\tag{L1}
\]
\end{lemma}
\begin{proof}
Standard; integrate the Lipschitz condition on the gradient. \hfill $\square$
\end{proof}

\begin{lemma}[Optimality condition of the proximal operator]\label{lem:prox_opt}
For any $v\in\R^{d}$ and $\eta>0$, $u=\prox_{\eta g}(v)$ iff
\[
\frac{1}{\eta}(v-u)\in\partial g(u).
\tag{L2}
\]
\end{lemma}
\begin{proof}
By definition of the proximal map as the minimiser of a strongly convex function. \hfill $\square$
\end{proof}

\begin{lemma}[Firm non‑expansiveness]\label{lem:firm}
For any $u,v\in\R^{d}$,
\[
\|\prox_{\eta g}(u)-\prox_{\eta g}(v)\|^{2}
\le\langle u-v,\;\prox_{\eta g}(u)-\prox_{\eta g}(v)\rangle .
\tag{L3}
\]
\end{lemma}
\begin{proof}
Follows from Lemma~\ref{lem:prox_opt} and monotonicity of $\partial g$. \hfill $\square$
\end{proof}

\section*{Main Proof (Step‑by‑Step Derivation)}
Let $x^{\star}$ be any minimiser of $F$. Denote $v_{k}=x_{k}-\eta\grad f(x_{k})$ and $x_{k+1}= \prox_{\eta g}(v_{k})$.
We prove the descent inequality and then telescope it.

\begin{enumerate}[label=[Step \arabic*],leftmargin=*,itemsep=1ex]
\item \textbf{Optimality condition of the proximal step.}\\
From Lemma~\ref{lem:prox_opt} applied to $v_{k}$ we have
\[
\frac{1}{\eta}\bigl(v_{k}-x_{k+1}\bigr)\in\partial g(x_{k+1}).
\tag{3}
\]

\item \textbf{Subgradient inequality for $g$.}\\
For any $x$, convexity of $g$ yields
\[
g(x)\ge g(x_{k+1})+\Big\langle \frac{1}{\eta}(v_{k}-x_{k+1}),\,x-x_{k+1}\Big\rangle .
\tag{4}
\]

\item \textbf{Descent lemma for $f$ at $x_{k+1}$.}\\
Using Lemma~\ref{lem:descent} with $x=x_{k}$ and $y=x_{k+1}$,
\[
f(x_{k+1})\le f(x_{k})+\langle\grad f(x_{k}),x_{k+1}-x_{k}\rangle
+\frac{L}{2}\|x_{k+1}-x_{k}\|^{2}.
\tag{5}
\]

\item \textbf{Combine \eqref{4} and \eqref{5}.}\\
Add the two inequalities (4) and (5):
\[
\begin{aligned}
F(x_{k+1}) 
&= f(x_{k+1})+g(x_{k+1})\\
&\le f(x_{k})+g(x_{k+1})
+\langle\grad f(x_{k}),x_{k+1}-x_{k}\rangle
+\frac{L}{2}\|x_{k+1}-x_{k}\|^{2}\\
&\le f(x_{k})+g(x)
+\Big\langle \frac{1}{\eta}(v_{k}-x_{k+1}),\,x-x_{k+1}\Big\rangle
+\langle\grad f(x_{k}),x_{k+1}-x_{k}\rangle
+\frac{L}{2}\|x_{k+1}-x_{k}\|^{2}.
\end{aligned}
\tag{6}
\]

\item \textbf{Replace $v_{k}=x_{k}-\eta\grad f(x_{k})$ and simplify inner products.}\\
Observe that
\[
\frac{1}{\eta}(v_{k}-x_{k+1})
= \frac{1}{\eta}\bigl(x_{k}-\eta\grad f(x_{k})-x_{k+1}\bigr)
= \frac{x_{k}-x_{k+1}}{\eta}-\grad f(x_{k}).
\]
Insert this into the inner product term of \eqref{6}:
\[
\Big\langle \frac{x_{k}-x_{k+1}}{\eta}-\grad f(x_{k}),\,x-x_{k+1}\Big\rangle
= \frac{1}{\eta}\langle x_{k}-x_{k+1},x-x_{k+1}\rangle
-\langle\grad f(x_{k}),x-x_{k+1}\rangle .
\tag{7}
\]

\item \textbf{Cancel the gradient terms.}\\
The term $-\langle\grad f(x_{k}),x-x_{k+1}\rangle$ from \eqref{7} together with $+\langle\grad f(x_{k}),x_{k+1}-x_{k}\rangle$ (the third term of \eqref{6}) gives
\[
-\langle\grad f(x_{k}),x-x_{k+1}\rangle
+\langle\grad f(x_{k}),x_{k+1}-x_{k}\rangle
= -\langle\grad f(x_{k}),x-x_{k}\rangle .
\tag{8}
\]

\item \textbf{Gather everything.}\\
Substituting \eqref{7}–\eqref{8} into \eqref{6} yields
\[
\begin{aligned}
F(x_{k+1}) &\le f(x_{k})+g(x)
+\frac{1}{\eta}\langle x_{k}-x_{k+1},x-x_{k+1}\rangle
-\langle\grad f(x_{k}),x-x_{k}\rangle
+\frac{L}{2}\|x_{k+1}-x_{k}\|^{2}.
\end{aligned}
\tag{9}
\]

\item \textbf{Use convexity of $f$.}\\
Convexity of $f$ gives
\[
f(x)\ge f(x_{k})+\langle\grad f(x_{k}),x-x_{k}\rangle .
\tag{10}
\]
Rearranging \eqref{10} and inserting into \eqref{9} eliminates $f(x_{k})$ and the gradient inner product:
\[
F(x_{k+1})\le F(x)+\frac{1}{\eta}\langle x_{k}-x_{k+1},x-x_{k+1}\rangle
+\frac{L}{2}\|x_{k+1}-x_{k}\|^{2}.
\tag{11}
\]

\item \textbf{Expand the inner product.}\\
Note that
\[
\langle x_{k}-x_{k+1},x-x_{k+1}\rangle
= \tfrac12\Bigl(\|x-x_{k}\|^{2}-\|x-x_{k+1}\|^{2}
-\|x_{k+1}-x_{k}\|^{2}\Bigr).
\tag{12}
\]

\item \textbf{Insert (12) into (11).}\\
\[
\begin{aligned}
F(x_{k+1}) &\le F(x)
+\frac{1}{2\eta}\Bigl(\|x-x_{k}\|^{2}-\|x-x_{k+1}\|^{2}
-\|x_{k+1}-x_{k}\|^{2}\Bigr)
+\frac{L}{2}\|x_{k+1}-x_{k}\|^{2}\\[0.3em]
&= F(x)+\frac{1}{2\eta}\bigl(\|x-x_{k}\|^{2}-\|x-x_{k+1}\|^{2}\bigr)
-\frac{1}{2\eta}\bigl(1-L\eta\bigr)\|x_{k+1}-x_{k}\|^{2}.
\end{aligned}
\tag{13}
\]

\item \textbf{Choose $x=x^{\star}$ and use $1-L\eta\ge0$ (Assumption A3).}\\
Since the last term is non‑positive, we obtain the fundamental inequality
\[
F(x_{k+1})-F(x^{\star})
\le \frac{1}{2\eta}\Bigl(\|x^{\star}-x_{k}\|^{2}
-\|x^{\star}-x_{k+1}\|^{2}\Bigr).
\tag{14}
\]

\item \textbf{Telescoping sum.}\\
Sum \eqref{14} for $k=0,\dots,K-1$:
\[
\sum_{k=0}^{K-1}\bigl(F(x_{k+1})-F(x^{\star})\bigr)
\le \frac{1}{2\eta}\bigl(\|x^{\star}-x_{0}\|^{2}
-\|x^{\star}-x_{K}\|^{2}\bigr)
\le \frac{\|x_{0}-x^{\star}\|^{2}}{2\eta}.
\tag{15}
\]

\item \textbf{Derive the $O(1/K)$ bound.}\\
Since the left‑hand side is the sum of $K$ non‑negative terms,
\[
K\;\min_{0\le k\le K-1}\bigl\{F(x_{k+1})-F(x^{\star})\bigr\}
\le \frac{\|x_{0}-x^{\star}\|^{2}}{2\eta},
\]
which implies
\[
\min_{0\le k\le K-1}\bigl\{F(x_{k+1})-F(x^{\star})\bigr\}
\le \frac{\|x_{0}-x^{\star}\|^{2}}{2\eta\,K}.
\tag{16}
\]
Choosing $k=K$ gives the explicit bound (2) of Theorem~\ref{thm:main}.
\hfill $\square$
\end{enumerate}

\section*{Rate Derivation (With Constants)}
The constant appearing in \eqref{2} is $C:=\|x_{0}-x^{\star}\|^{2}/(2\eta)$.  
If the stepsize is taken at the maximal admissible value $\eta=1/L$, then
\[
F(x_{k})-F(x^{\star})\le \frac{L\,\|x_{0}-x^{\star}\|^{2}}{2k},
\]
showing that a larger $L$ (steeper curvature) slows the convergence, exactly as in gradient descent.

\section*{Special Cases}
\begin{enumerate}[label=\arabic*.]
\item \textbf{Pure gradient descent ($g\equiv0$).}  
Equation \eqref{14} reduces to the classic smooth convex descent inequality, yielding the same $O(1/k)$ rate.
\item \textbf{Indicator function of a convex set $C$.}  
If $g=\iota_{C}$, then $\prox_{\eta g}$ is the Euclidean projection onto $C$, and the method becomes the projected gradient algorithm with the same convergence guarantee.
\item \textbf{Lasso ($g(x)=\lambda\|x\|_{1}$).}  
The proximal map is soft‑thresholding, which is cheap to compute; the bound \eqref{2} holds unchanged.
\end{enumerate}

\end{document}